\section{Q\&A：真正困难的是评审、注意力与安全}

主讲结束后的问答没有新增教学 slide，却补齐了系统部署最重要的约束：怎样证明 novelty，Agent 如何共享知识，科学家怎样阅读百页报告，计算规模为何未公开，以及多智能体系统怎样阻止危险研究目标。这里的口头说明是工程经验，不是形式化保证。

\subsection{Knowledge cutoff 评测并不自动干净}

一种自然想法是只让模型读取某个历史年份之前的材料，再测试它能否预测之后的发现。但模型训练语料、后续微调、检索缓存、论文引用和公开讨论都可能泄漏未来信息；同时，真正重要的发现往往依赖当时尚不存在的实验工具。时间切分有价值，但必须冻结整个系统而不是只写一句“使用旧知识”。

\begin{warningbox}{时间切分的最低要求}
冻结基础模型、检索索引、工具版本和输入文档；记录运行时间；选择在 cutoff 后才公开且此前没有摘要泄漏的目标；让评审看不到目标答案；保留所有候选而非只展示命中结果。否则所谓 prediction after cutoff 可能只是隐性 retrieval。
\end{warningbox}

\teachervoice{01:00:02--01:01:18，老师认为 knowledge cutoff 是有吸引力的评测方向，同时提醒模型更新和公开痕迹让真正的 temporal holdout 很难。}

\subsection{Agent communication 与 peer-review bandwidth}

当多个 Agent 并行产生大量假设，通信问题会从“怎样发消息”升级为“怎样建立共享但不过度同质化的研究记忆”。若所有 Agent 只读同一 summary，早期错误会扩散；若完全隔离，又会重复搜索。Peer review 也会成为瓶颈：人类不可能逐页审查无限增长的自动报告。

\begin{knowledgebox}{共享 memory 的三层设计}
事实层保存带来源的观察；假设层保存可竞争的机制和反例；过程层保存谁在何时用什么模型、工具与预算得到结论。Agent 可以共享事实而保留独立假设，再由 meta-review 比较分歧，避免一份统一摘要过早消灭探索多样性。
\end{knowledgebox}

\teachervoice{01:01:18--01:02:49，老师指出未来可能需要面向 Agent 的 peer-review system 与 knowledge base，因为自动生成速度会超过人类审稿速度。}

\subsection{报告排序既要指导注意力，也要保留弱信号}

科学家需要系统指出最值得先读的部分，但低排名候选可能含有尚未被 judge 理解的关键线索。合理界面应同时提供 top recommendation、反对理由、低排名 novelty bucket 与停止建议。系统也应允许回答“目前没有足够好想法，请重新定义问题”。

\begin{lstlisting}[language=Python,caption={兼顾优先级与弱信号的报告分层}]
def triage(hypotheses):
    return {
        "act_now": select_high_evidence_high_value(hypotheses),
        "test_cheaply": select_high_information_low_cost(hypotheses),
        "novel_but_uncertain": preserve_outliers(hypotheses),
        "rejected_with_reasons": keep_failure_provenance(hypotheses),
        "reframe_goal": detect_no_viable_path(hypotheses),
    }
\end{lstlisting}

\teachervoice{01:02:53--01:04:30，老师说报告会告诉科学家先看哪些路径，但仍保留被系统判为不可行的想法，因为其中可能有线索；数学问题有时需要人把大问题拆成恰当子问题，系统才能推进。}

\begin{importantbox}{Human decomposition 仍是核心能力}
Agent 能扩展搜索，却未必知道什么是正确子问题。专家通过改变边界条件、选择可测代理、拆分证明或限定生物系统，决定搜索空间是否可验证。人机协作的关键不是人负责“最后确认”，而是人持续塑造问题表示。
\end{importantbox}

\subsection{没有公开的 Agent 数量，不能靠猜}

听众追问一次 overnight run 使用多少并行 Agent、多少 token。主讲人表示在录制结束后再谈，没有公开数字。这个空白本身很重要：没有并行度、token、工具调用和成本，就无法判断 quality gain 来自架构还是单纯增加计算，也无法复现实验经济性。

\begin{warningbox}{Compute transparency}
科研 Agent 报告至少应公开模型版本、并行 worker、总 token、墙钟时间、工具调用、失败重试和人工时。若商业或安全原因不能公开精确值，也应提供区间和归一化效率。讲义不根据“overnight”反推虚构数字。
\end{warningbox}

\teachervoice{01:04:34--01:04:54，老师明确没有在录制中披露 agent 数与 token 量。这里必须保留为 source gap。}

\subsection{多层安全：从目标过滤到运行中监控}

开放式科学系统可能被用于病原体、毒物或其他危险研究。课程描述三层防线：输入目标 screening；运行中持续监控生成内容；当危险想法比例超过报告中的阈值时停止计算；此外继承 base Gemini 的安全训练与测试。讲者也承认，多 Agent 和工具扩大了出错与滥用表面。

\begin{lstlisting}[language=Python,caption={课堂所述多层安全监控的概念化版本}]
def run_research(goal):
    if goal_filter.unsafe(goal):
        return halt("unsafe research goal")

    unsafe, total = 0, 0
    for idea in co_scientist.stream(goal):
        total += 1
        unsafe += int(content_monitor.unsafe(idea))
        if total >= minimum_window and unsafe / total >= 0.10:
            return halt("unsafe trajectory")
        yield sandbox_and_log(idea)
\end{lstlisting}

\begin{warningbox}{10\% 阈值不是安全定律}
课堂提到的阈值是特定实现细节。危险内容可能在比例很低时就造成严重后果，也可能通过分散、编码或工具链绕过分类器。真正的 defense in depth 还需要权限隔离、工具沙箱、敏感数据控制、人工审批、红队、审计和事后响应。
\end{warningbox}

\teachervoice{01:04:56--01:06:20，老师说明输入检查可能被改写绕过，所以还要监控运行中的 ideas；同时明确承认 multi-agent setup 扩大了可能出错和被滥用的表面。}

\subsection{本章小结}

Q\&A 把问题从“Agent 会不会生成好想法”推进到部署现实：时间切分可能泄漏，Agent communication 可能放大共识错误，peer review 与人类注意力会成为瓶颈，compute 信息可能不完整，安全过滤也会被绕过。可靠系统必须把 provenance、资源、负结果、报告分层和安全控制设计为一等对象。

